\section{The reduction}\label{sec:reduction}
Fix any decidable predicate $R(e,n,t)$ and the set $S=\{e:\exists n\,\forall t\ R(e,n,t)\}$. We effectively choose the cellular automaton in Section~\ref{sec:tree} to run the following program, with $e$ hardwired.

\paragraph{Read a positive-time return.}
At time $y$ a waiting front is at position $y$. Ignore the time-zero pulse. At the first positive-time pulse, freeze the front and retain its coordinate as unary input $q$. With the transition convention above the mark is written into row $q+1$ at position $q$. Test $q$. If $q=1$ or $q$ is composite, halt. Otherwise compute the unique $n\geq0$ with $q=p_{n+1}$, the $(n+1)$st prime. Then evaluate $R(e,n,0),R(e,n,1),\ldots$ and halt at the first false answer. All later computation states ignore the external pulse.

For a $\Sigma^0_2$-complete choice of $S$, there is no computable bound on a sufficient clock order as a function
of the input game: such a bound would leave only finitely many possible
$n$ in this reduction, making $\exists n\,\forall t\ R(e,n,t)$ a
$\Pi^0_1$ predicate. Indeed one tests, for every finite time cutoff, whether
some one of those finitely many candidates survives all checks so far;
if none survives forever, all have failed by a common finite cutoff.
This contradicts $\Sigma^0_2$-hardness.

This finite program has an effective finite radius-one simulation. The initial-front rule, frozen unary track and ordinary tape/head tracks give the stated boundary conventions. We do not require a bound on its running time. Each completed false check causes the halt symbol after finitely many updates.

Include the Fock, coordinate, clock, reader and tableau equations above, all typing relations, and the stated $Z_{\rm clk}$ commutators. They form an admissible finite presentation. For example their transported polynomials include
\[
 W_4a_{ab}=a_{ab}Z_{{\rm clk},0}^bW_0^2,\qquad
 C_{{\rm ctr},4}a_{ab}=2^{-b}a_{ab}C_{{\rm ctr},0}^2,\qquad
 J_{{\rm W},4}a_{ab}=a_{ab}J_{{\rm W},0}.
\]
The source-only return equations and tableau equations are centralizer polynomials. Lemma~\ref{lem:four} therefore computes a finite rational Boolean game $G_e$.

\begin{theorem}\label{thm:reduction}
The map $e\mapsto G_e$ satisfies
\[
 G_e\in\perf_{\qs}\ \Longleftrightarrow\
 G_e\in\perf_{\finS}\ \Longleftrightarrow\
 \exists n\,\forall t\ R(e,n,t).
\]
No $G_e$ has a perfect finite-dimensional strategy.
\end{theorem}
\begin{proof}
A perfect normal strategy extracts a nonzero faithful root realization. The determinant argument supplies a nonzero root phase sector of finite order $q$. By \eqref{eq:weyl} its word descendants persist at all depths. Equations~\eqref{eq:W}--\eqref{eq:return} force the exact return sequence, with its first positive return at $q$. Lemma~\ref{lem:tableau} needs only that the initial root vector lie in this $Z_{\rm clk}$-eigenspace: the clock pulse is then scalar on all its word descendants. The imposed $[Z_{\rm clk},P_c]=0$ relations are redundant for this induction. An invalid $q$ halts. For valid $q=p_{n+1}$ any false $R(e,n,t)$ also halts. A sufficiently large square would contain that halt, contrary to $P_H=0$. Hence the extracted $n$ witnesses membership in $S$.

Conversely take a witness $n$ and set $q=p_{n+1}$. Use \eqref{eq:fockdensity} with $\sigma=I_q/q$, a good primitive phase and its Weyl pair, the canonical coordinate permutations, \eqref{eq:coordinateclock}, \eqref{eq:return}, \eqref{eq:D}, and the classical nonhalting tableau. Every required operator is bounded and commutes with the density when so specified. This is a simultaneous realization. On $\mathbb C^q$ one may take $J_{\rm W}|j\rangle=\zeta^j|j\rangle$, $K_{\rm W}|j\rangle=|j+1\bmod q\rangle$, and $Z_{\rm clk}=\zeta I$. Its root entropy and largest Schmidt probability are
\begin{equation}\label{eq:witnessentropy}
 S(h)=\log_2q+\log_2(4/3)+4/3,\qquad p_{\max}=3/(4q).
\end{equation}
The compiler adds finite factors, giving a perfect finite-entropy strategy.

The tensor-product witness can be written without an implicit extension.
Let $\mathcal H_F=\ell^2(\mathcal W)$ and let $\mathcal H_{\rm sc}=\mathbb C^{17}$
index the scaffold corners. Take
\[
 \mathcal H_A=\mathcal H_B=
 \mathcal H_{\rm sc}\otimes\mathcal H_F\otimes\mathbb C^q
 \otimes\mathbb C^3\otimes\mathbb C^{d_{\max}}.
\]
If $\mu=(1,2,2,2,2,4,4,4,4,8,8,8,8,16,16,16,16)/121$,
use the pure vector
\[
 |\psi\rangle=|\sqrt\mu\rangle\!\rangle\otimes
 |\sqrt{h_F}\rangle\!\rangle\otimes|\Phi_q\rangle
 \otimes|\Phi_3\rangle\otimes|\Phi_{d_{\max}}\rangle,
 \quad |\sqrt\mu\rangle\!\rangle=\sum_i\sqrt{\mu_i}|i,i\rangle.
\]
The Fock factor is $\sum_w\sqrt{(3/4)16^{-|w|}}|w,w\rangle$.
Alice uses the scaffold FMS effects, actual graph projections, refinements
and column partitions of Lemma~\ref{lem:four}; Bob uses the prescribed
opposite FMS effects and the transposed transferred centralizer products.
The row factor is the finite maximally entangled ancilla of
Lemma~\ref{lem:rows}. These simultaneous effects were checked in the
compiler proof. The marginal entropy is
$H(\mu)+S(h_F)+\log_2q+\log_2 3+\log_2d_{\max}<\infty$.

Finally finite-dimensional perfect play would extract four isometries with mutually orthogonal ranges on one nonzero finite-dimensional root space. Their range dimensions would sum to four times the dimension of that space, which is impossible. The argument applies to every input, whether or not it has a witness.
\end{proof}

Together with Section~\ref{sec:upper}, this proves both completeness statements in Theorem~A. The value-one promise is established next. For the promise version, the reduction maps every input into that promise; the unrestricted $\Sigma^0_2$ membership predicate restricts to it, giving promise-problem completeness. The construction can be made syntactically nonsynchronous and non-imitation on every input by adding a fresh binary question to each side, accepting every pair involving it. Either answer then has the same nonempty allowed-answer support against every opposite question, which violates the imitation separation condition. This addition does not change perfection or the value-one property.

\paragraph{Size of the finite game.}
Here is a coarse order-of-magnitude bound, independent of the witness $q$.
Let $g=|\Gamma|$, let $N_e$ count the named graph projections, rows and
word prefixes after endpoint translation, and let $b_e$ bound the bit
length of their rational coefficients. The lookup table has at most
$2(g+1)g^2$ entries, so unshared explicit polynomials give
$N_e=\operatorname{poly}(g)$ for this fixed list of operator relations.
Denominator clearing and signed repetitions give row dimension
$d_{\max}\leq2^{O(N_e(b_e+\log(N_e+1)))}$.
The scalar contexts, column questions and marginal bindings are polynomial
in $N_e,d_{\max}$, so the padded question count $Q_e$ and maximum answer
count obey the same coarse exponential bound. An explicit predicate table
is polynomial in those two counts and has the same exponential order.
Thus the effective compilation depends on the hardwired program, never on
its chosen prime or running time. For scale, the bundled minimal
one-isometry example already has $748$ rows and $27\,991$ questions per
party; it is not a size claim for the full reduction.
